% !PRINTED 62
\weberpage{62}

Wir fragen zunächst nach der Möglichkeit von $T$-Funktionen nullter Ordnung: Ist $L(u)$ eine solche und $L'(u)$ ihre Ableitung, so folgt durch zweimalige Differentiation von (6)

\[
\frac{d}{du}\,\frac{L'(u+\omega_1)}{L(u+\omega_1)} = \frac{d}{du}\,\frac{L'(u)}{L(u)},
\]

und entsprechend aus der zweiten Gleichung (6). Danach ist $d^2\log L(u)/du^2$ als ganze doppeltperiodische Funktion eine Konstante, und folglich $\log L(u)$ eine ganze Funktion zweiten Grades von $u$. Wir erhalten also:

4. \gesperrt{Eine $T$-Funktion nullter Ordnung ist von der Form}

\weq{7}{L(u) = C\,e^{-\pi i (\lambda u^2 + \mu u)},}

worin $\lambda$, $\mu$, $C$ Konstanten sind.

Daß jede Funktion dieser Form eine $T$-Funktion nullter Ordnung ist, ersieht man ohne weiteres.

Das Produkt

\weq{8}{T'(u) = L(u)\,T(u)}

ist ebenso wie $T(u)$ eine $T$-Funktion $m$ter Ordnung und verschwindet in denselben Punkten.

Die Periodizitätsfaktoren von $T'$ erhält man aus denen von $T$, wenn man die Exponenten $a_1$, $a_2$, $b_1$, $b_2$ durch

\weq{9}{\begin{aligned}
a_1' &= a_1 + \lambda\omega_1, & b_1' &= b_1 + \mu\omega_1\\
a_2' &= a_2 + \lambda\omega_2, & b_2' &= b_2 + \mu\omega_2
\end{aligned}}

ersetzt, worin die $b$, $b'$ jedoch nur bis auf additive gerade ganze Zahlen bestimmt sind.

5. \gesperrt{Nach (9) kann man $\mu$ immer und nur auf eine Weise so bestimmen, daß $b_1'$, $b_2'$ reell werden.}

Denn setzt man die imaginären Teile von $b_1'$, $b_2'$ gleich Null, so erhält man für den imaginären und reellen Teil von $\mu$ zwei lineare Gleichungen, deren Determinante nach \S\,16 der Flächeninhalt des Periodenparallelogramms, also von Null verschieden ist.

Die Periodizitätsfaktoren und die Nullpunkte einer $T$-Funktion $m$ter Ordnung stehen in einer gewissen Abhängigkeit voneinander, die wir leicht aus den Sätzen des \S\,15 erhalten. Zunächst ergibt sich die Ordnung $m$, wenn man das Integral

\[
\frac{1}{2\pi i}\int d\log T(u) = m
\]
%CONT>>
